% Kapitel 1 -- Von der Hexe zu Ottenstein und wie Kuonr\^{a}t zum Herrn von Raabs ward

\section{Von der Hexe zu Ottenstein und Kuonr\^{a}ts Ritterschlag}

\mylettrine{N}{un} ward \gls{Kuonrat von Steinare} von \gls{Friedrich II von Oesterreich} selbst nach Burg Ottenstein zurückgebracht, wie mir glaubwürdige Quellen berichteten. Bei seiner Ankunft wurde vernommen, daß im selben Jahre des Herrn 1232 der Herr \gls{Wolfhart von Raabs} verschieden sei. Der Herzog ehrte das Edikt seines Vaters und schlug den jungen Mann zum Ritter; zugleich erhob er ihn zum Herrn von \gls{Raabs an der Thaya}. Der Termin der Vermählung ward auf zwei Monde bestimmt; doch gebot der Herzog, zum Verdruß des jungen Herrn, daß dieser seine Braut vor der Hochzeit nicht sehen dürfe, und er solle bis dahin bei seiner Familie auf Ottenstein verbleiben.

\gls{Kuonrat von Steinare} gehorchte, murrend, denn seine Gedanken waren ohnehin anderswohin gewandt. Der \gls{Knecht} hatte während der Gefangenschaft seines Herrn die Zeit auf Ottenstein im Stall zugebracht; gleichwohl hielt ihm der Herr vor, nichts getan zu haben. Auch sein Ross \gls{Veltloufer} litt unter lahmen Hufen, und Kuonr\^{a}t schrieb dies dem wilden Ritt von \gls{Znaim} nach Ottenstein zu, als er vor die Botschaft von Kuonr\^{a}ts Gefangenschaft überbrachte.

So blieb es nicht bei mahnenden Worten. In seinem Zorn ergriff \gls{Kuonrat von Steinare} den \gls{Knecht} und züchtigte ihn. Ob die Schläge nicht mehr als eine grobe Lehre waren, wird verschiedentlich erzählt; so heißt es, der Knecht habe die Demut bewahrt und sich nicht gerächt.

Während der langen Wartezeit aber soll \gls{Kuonrat von Steinare} von einer Unruhe geplagt worden sein, die weder im Gebet noch in der Arbeit Linderung fand. Man weiß wohl, was einen Mann seines Wesens und seines jungen Blutes zu plagen pflegt. So suchte er die \gls{Schenke} zu Ottenstein auf, die an jenem Abend sehr voll gewesen sein soll, denn die Tagelöhner der Gegend tranken ihren Lohn.

\gls{Kuonrat von Steinare} trank, wie mir ein zuverlässiger Zeuge berichtete, rascher als Klugheit riet. Es ist nämlich bekannt, daß er dem Trunke wenig gewachsen war; schon der zweite Becher ließ sein Urteil schwinden. Im Rausch vergaß er, daß eine Braut auf ihn wartete, und erklärte laut, es sei seine \gls{Freiheit}, auch als Verlobter nach Manneslust zu handeln. Er rief nach einer Dirne und forderte, sie solle mit ihm gehen. Die Frau hingegen verlangte Bezahlung für ihren Dienst, was freilich zu begreifen ist, wenngleich solches Treiben des Lobes nicht bedarf.

Dies empfand \gls{Kuonrat von Steinare} als bittere Beleidigung. Er beschimpfte sie als Hexe; sie habe ihn mit einem Wohlgeruch bezaubert, damit er ihr verfalle, und er sei ihr darum nichts schuldig. Mit diesem Schrei verließ er die \gls{Schenke}, doch nicht ohne Beschwernis: er stolperte über drei Hocker, fiel polternd zur Seite und stieß dabei einen Musikanten, der dort seine Laute angestimmt hatte, unversehens in das Herdfeuer. Der Glückliche verbrannte sich, dem Herrn sei Dank, nur leicht am Arm und war nach wenigen Tagen wieder heil. \gls{Kuonrat von Steinare} indes war bereits verschwunden, ohne sich zu besinnen.

Am folgenden Morgen war die Geschichte in Ottenstein allseits bekannt. \gls{Kuonrat von Steinare} hielt unbeirrt daran fest, die Frau habe ihn bezaubert; er bezeichnete sie öffentlich als Hexe und behauptete, sie habe ihn durch einen teuflischen Wohlgeruch ihrer Macht untertan gemacht. So eilte er zum Burgkaplan, auf daß dieser gegen die Hexe vorgehe. Doch der Kaplan war nicht zugegen.

Stattdessen traf er den jungen Priester \gls{Kuono}, der auf Durchreise war und sich auf dem Weg nach \gls{Retz} befand. Dieser hörte die Klage mit stiller Aufmerksamkeit an und verstand, was wirklich geschehen war, so heißt es jedenfalls unter denen, die ihn kannten. Er versicherte \gls{Kuonrat von Steinare}, die Sache zu regeln. Dem Kaplan gegenüber erwirkte er, daß die Frau förmlich der Hexerei angeklagt werde; zur Frau selbst aber begab er sich in aller Stille, gab ihr ein Pferd und ein Nonnengewand und wies ihr den Weg nach \gls{Retz}. Sie erkannte die Gefahr und tat, was Klugheit gebot: sie verließ Ottenstein noch in derselben Nacht.

Als \gls{Kuonrat von Steinare} nach dem Verbleib der Hexe fragte, antwortete ihm \gls{Kuono} mit ernstem Gesicht, die Frau sei auf ihrem Besen davongeflogen und niemand habe sie mehr gesehen. Der junge Herr glaubte ihm auf der Stelle.

So rettete \gls{Kuono} die Unschuldige und bewahrte zugleich die Ehre seines jungen Gefährten. Daß er dabei den Aberglauben des neuen Herrn von \gls{Raabs an der Thaya} nur noch tiefer einpflanzte, mag der Herr ihm nachsehen; der Nutzen überwog das Übel in jenem Augenblick.
